\documentclass[a4paper,11pt]{article}
\usepackage[T2A]{fontenc}
\usepackage[utf8]{inputenc}
\usepackage[russian]{babel}
\usepackage{amsmath,amssymb,mathtools}
\renewcommand{\familydefault}{\sfdefault}
\usepackage{geometry}
\geometry{left=22mm,right=20mm,top=22mm,bottom=23mm}
\usepackage{microtype}
\usepackage{icomma}
\usepackage{booktabs,tabularx,array}
\usepackage{enumitem}
\usepackage{xcolor}
\usepackage{tikz}
\usepackage{hyperref}
\usepackage{fancyhdr}
\usepackage{setspace}
\setstretch{1.14}
\definecolor{accent}{HTML}{006F75}
\definecolor{darkaccent}{HTML}{00464D}
\definecolor{ink}{HTML}{25343A}
\definecolor{err}{HTML}{AD2C32}
\hypersetup{colorlinks=true,linkcolor=darkaccent,pdfauthor={Лёвин Артём Александрович},pdftitle={Ревью домашней работы — Ярослав Гаврилов — 15.11.2026}}
\setlength{\parindent}{0pt}
\setlength{\parskip}{2pt}
\setlist[itemize]{nosep,topsep=3pt,leftmargin=1.5em}
\pagestyle{fancy}
\fancyhf{}
\fancyhead[L]{\small\color{darkaccent}Ревью домашней работы}
\fancyhead[R]{\small\color{ink}15.11.2026}
\fancyfoot[C]{\small\color{ink}\thepage}
\renewcommand{\headrulewidth}{0.25pt}
\setlength{\headheight}{15pt}
\newcommand{\rlabel}[1]{\par\noindent\textbf{\color{darkaccent}#1}\quad}
\newcommand{\errorlabel}{\textcolor{err}{\textbf{Ошибка / недочёт:}}}
\newcommand{\keylabel}{\textcolor{blue!60!black}{\textbf{Ключевая идея:}}}
\newcommand{\finallabel}{\textcolor{teal!60!black}{\textbf{Итог:}}}
\newcommand{\tasktitle}[2]{\clearpage\phantomsection\label{task:#1}\section*{Задание №\,#1. #2}\vspace{-3mm}\textcolor{accent}{\rule{\linewidth}{0.45pt}}}
\begin{document}
\thispagestyle{empty}
\begin{center}
\vspace*{30mm}
{\Large\bfseries\color{darkaccent}МАТЕМАТИКА И ФИЗИКА}\par
\vspace{24mm}
{\huge\bfseries Ревью домашней работы}\par
\vspace{6mm}
{\Large Ярослав Гаврилов}\par
\vspace{22mm}
{\large 15 ноября 2026 года}\par
\vfill
{\large Лёвин Артём Александрович}\par
{\small эксклюзивно для Ярослава Гаврилова}\par
\vspace{12mm}
\begin{tikzpicture}
\draw[accent!65,line width=0.9pt] (-5,0) .. controls (-2.6,0.7) and (-1.1,-0.55) .. (0,0) .. controls (1.6,0.7) and (3.0,-0.5) .. (5,0);
\end{tikzpicture}
\end{center}
\clearpage
\section*{Сводная таблица}
В таблице показаны только задания, требующие исправления или завершения. Номера ведут к подробному разбору.\par
\vspace{5mm}\noindent
{\small
\setlength{\tabcolsep}{2.2pt}
\renewcommand{\arraystretch}{1.28}
\begin{tabular}{@{}>{\centering\arraybackslash}p{9mm} >{\raggedright\arraybackslash}p{31mm} >{\raggedright\arraybackslash}p{58mm} >{\raggedright\arraybackslash}p{28mm} >{\raggedright\arraybackslash}p{26mm}@{}}
\toprule
\textbf{№} & \textbf{Тема} & \textbf{Суть ошибки / недочёта} & \textbf{Ваш ответ} & \textbf{Верный ответ}\\
\midrule
\hyperref[task:2]{2} & Тригонометрия и корни & Вычисление не доведено до результата & Не указан & $-2$\\
\hyperref[task:4]{4} & Двойной угол & При раскрытии скобок потерян множитель $30$ & Не указан & $-27{,}6$\\
\hyperref[task:8]{8} & Кинетическая энергия & Скорость записана, энергия не найдена & Не указан & $0{,}242$ Дж\\
\hyperref[task:9]{9} & Степенной закон газа & Решение не представлено & Нет решения & $8$ м$^3$\\
\hyperref[task:11]{11} & Степенной закон газа & Решение не представлено & Нет решения & $5$\\
\hyperref[task:12]{12} & Дифракция & Не обоснован выбор периода и не найден угол & Не указан & $30^\circ$\\
\hyperref[task:13]{13} & Показательная функция & Неверно найден параметр $b$; решение не закончено & Не указан & $28$\\
\hyperref[task:14]{14} & График тангенса & Решение не представлено & Нет решения & $1$\\
\bottomrule
\end{tabular}
}
\clearpage
\thispagestyle{empty}
\vspace*{0.37\textheight}
\begin{center}
{\Huge\bfseries\color{darkaccent}Разбор ошибок}
\end{center}

\tasktitle{2}{Тригонометрическая функция и радикалы}
\rlabel{Текст задания.} Найдите значение выражения
\[
\sqrt8-\sqrt{32}\sin^2\!\frac{11\pi}{8}.
\]
\rlabel{Ваше решение.} Вы правильно применили формулу понижения степени
\[
\sin^2u=\frac{1-\cos2u}{2}
\]
и дошли до преобразования
\[
\sqrt8-\sqrt{32}\cdot\frac{1+\frac{\sqrt2}{2}}{2}.
\]
\rlabel{Ваш ответ.} Не указан.
\par\noindent\errorlabel\ Вычисление остановлено перед упрощением квадратных корней и приведением подобных членов. В видимой части преобразований существенной математической ошибки нет, но задание не завершено.
\rlabel{Где именно остановились.} Последний правильный шаг --- представление квадрата синуса через косинус двойного угла и использование
\[
\cos\frac{11\pi}{4}=\cos\frac{3\pi}{4}=-\frac{\sqrt2}{2}.
\]
Первый недостающий шаг --- заменить $\sqrt8$ на $2\sqrt2$, а $\sqrt{32}$ на $4\sqrt2$.
\rlabel{Почему это важно.} Исходное выражение содержит несколько радикалов. Если оставить их без упрощения, не видно сокращения слагаемых с $\sqrt2$. Ответ здесь получается целым числом, однако это обнаруживается только после аккуратного раскрытия скобок.
\par\noindent\keylabel\ Сначала преобразуйте тригонометрическую часть, затем вынесите полные квадраты из-под корней и приведите подобные слагаемые.
\rlabel{Исправление от последнего правильного шага.}
\[
\begin{aligned}
\sqrt8-\sqrt{32}\cdot\frac{1+\frac{\sqrt2}{2}}2
&=2\sqrt2-4\sqrt2\cdot\frac{1+\frac{\sqrt2}{2}}2\\
&=2\sqrt2-2\sqrt2\left(1+\frac{\sqrt2}{2}\right).
\end{aligned}
\]
\rlabel{Полное правильное решение.} По формуле понижения степени
\[
\sin^2\frac{11\pi}{8}
=\frac{1-\cos\frac{11\pi}{4}}2
=\frac{1+\frac{\sqrt2}{2}}2.
\]
Тогда
\[
\begin{aligned}
\sqrt8-\sqrt{32}\sin^2\frac{11\pi}{8}
&=2\sqrt2-2\sqrt2\left(1+\frac{\sqrt2}{2}\right)\\
&=2\sqrt2-2\sqrt2-2=-2.
\end{aligned}
\]
\rlabel{Проверка.} $\sin^2(11\pi/8)=(2+\sqrt2)/4$; подстановка даёт то же выражение $2\sqrt2-\sqrt2(2+\sqrt2)=-2$.
\par\noindent\finallabel\ $\boxed{-2}$.
\rlabel{Задача на закрепление.} Вычислите $\displaystyle\sqrt{18}-\sqrt{72}\sin^2\frac{5\pi}{8}$.

\tasktitle{4}{Формула косинуса двойного угла}
\rlabel{Текст задания.} Найдите $30\cos2\alpha$, если $\cos\alpha=\frac15$.
\rlabel{Ваше решение.} Вы использовали тождество $\cos2\alpha=2\cos^2\alpha-1$, после чего в исправленном варианте записали $60\cos^2\alpha-1$, а затем $\frac{60}{25}-1$. Значение до конца не вычислено.
\rlabel{Ваш ответ.} Не указан.
\par\noindent\errorlabel\ Множитель $30$ распределён только на первое слагаемое внутри скобок.
\rlabel{Где именно возникает ошибка.} Последний правильный шаг:
\[
30\cos2\alpha=30\bigl(2\cos^2\alpha-1\bigr).
\]
Первый неправильный переход:
\[
30\bigl(2\cos^2\alpha-1\bigr)\ne60\cos^2\alpha-1.
\]
\rlabel{Почему это неверно.} Согласно распределительному свойству умножения, каждый член разности умножается на число перед скобками:
\[
c(u-v)=cu-cv.
\]
Поэтому $30\cdot(-1)=-30$, а не $-1$. Дальнейшее аккуратное приведение дробей не исправит ошибку, сделанную до него.
\par\noindent\keylabel\ Сначала раскройте скобки полностью, затем подставляйте данное значение косинуса.
\rlabel{Исправление от последнего правильного шага.}
\[
30\bigl(2\cos^2\alpha-1\bigr)=60\cos^2\alpha-30.
\]
\rlabel{Полное правильное решение.} Так как $\cos\alpha=\frac15$, то
\[
\cos^2\alpha=\left(\frac15\right)^2=\frac1{25}.
\]
По формуле двойного угла
\[
\cos2\alpha=2\cdot\frac1{25}-1
=\frac2{25}-\frac{25}{25}=-\frac{23}{25}.
\]
Следовательно,
\[
30\cos2\alpha=30\cdot\left(-\frac{23}{25}\right)
=-\frac{138}{5}=-27{,}6.
\]
\rlabel{Проверка.} $-1\le\cos2\alpha\le1$, поэтому $30\cos2\alpha$ находится от $-30$ до $30$. Полученное число удовлетворяет этому ограничению. Кроме того, $\cos^2\alpha=0{,}04<0{,}5$, поэтому косинус двойного угла действительно отрицателен.
\par\noindent\finallabel\ $\boxed{-27{,}6}$.
\rlabel{Задача на закрепление.} Найдите $18\cos2\beta$, если $\cos\beta=\frac13$.

\tasktitle{8}{Скорость колебаний и кинетическая энергия}
\rlabel{Текст задания.} Скорость груза массой $m=0{,}4$ кг изменяется по закону $v=v_0\cos\frac{2\pi t}{T}$, где $v_0=1{,}1$ м/с, $T=2$ с. Найдите кинетическую энергию $E=\frac{mv^2}{2}$ в момент $t=17$ с.
\rlabel{Ваше решение.} Вы выписали обе формулы и правильно получили
\[
v=1{,}1\cos\frac{2\pi\cdot17}{2}=1{,}1\cos17\pi.
\]
\rlabel{Ваш ответ.} Не указан.
\par\noindent\errorlabel\ Решение не доведено до значения скорости и вычисления энергии. Последний правильный шаг уже позволяет найти результат без дополнительных формул.
\rlabel{Где именно остановились.} После подстановки $t=17$ осталось вычислить $\cos17\pi$, затем подставить скорость в формулу $E=\frac{mv^2}{2}$.
\rlabel{Почему важно закончить.} Скорость может быть отрицательной, поскольку она указывает и направление движения. Кинетическая энергия, напротив, неотрицательна: в её формулу входит квадрат скорости. Нельзя принимать знак скорости за знак энергии.
\par\noindent\keylabel\ Для нечётного целого $n$ выполняется $\cos n\pi=-1$. Далее скорость возводится в квадрат.
\rlabel{Исправление от последнего правильного шага.}
\[
v=1{,}1\cos17\pi=-1{,}1\text{ м/с}.
\]
\rlabel{Полное правильное решение.} Находим аргумент косинуса:
\[
\frac{2\pi t}{T}=\frac{2\pi\cdot17}{2}=17\pi.
\]
Поскольку $17$ нечётно,
\[
v=1{,}1\cos17\pi=-1{,}1\text{ м/с}.
\]
Теперь вычисляем кинетическую энергию:
\[
E=\frac{0{,}4\cdot(-1{,}1)^2}{2}
=\frac{0{,}4\cdot1{,}21}{2}
=\frac{0{,}484}{2}=0{,}242\text{ Дж}.
\]
\rlabel{Проверка.} Максимальный модуль скорости равен $v_0=1{,}1$ м/с, поэтому энергия не может превышать $\frac{mv_0^2}{2}=0{,}242$ Дж. В момент $t=17$ с модуль скорости равен максимальному, и найденная энергия совпадает с этим значением.
\par\noindent\finallabel\ $\boxed{0{,}242\text{ Дж}}$.
\rlabel{Задача на закрепление.} Скорость груза массой $0{,}5$ кг задана формулой $v=2\cos\frac{2\pi t}{4}$ (м/с). Найдите его кинетическую энергию через $6$ с.

\tasktitle{9}{Степенной закон адиабатического процесса}
\rlabel{Текст задания.} Для газа выполняется $pV^k=3{,}2\cdot10^6$, где $k=\frac43$. Найдите $V$ в м$^3$, если $p=2\cdot10^5$ Па.
\rlabel{Ваше решение.} На представленных фотографиях решения этого задания нет.
\rlabel{Ваш ответ.} Не указан.
\par\noindent\errorlabel\ Отсутствует необходимый ход решения. Здесь требуется выразить степень неизвестного объёма, а затем избавиться от дробного показателя.
\rlabel{Где именно возникает недочёт.} Последнего правильного шага в работе нет. Первый необходимый шаг --- подставить давление в равенство $pV^{4/3}=3{,}2\cdot10^6$.
\rlabel{Почему нельзя сразу делить показатели.} В уравнении неизвестен объём $V$, а не показатель степени. Сначала надо изолировать $V^{4/3}$; затем обе части равенства возвести в обратную степень $\frac34$. Положительность объёма позволяет использовать эти преобразования без неоднозначности.
\par\noindent\keylabel\ Если $V^{4/3}=A$ и $V>0$, то $V=A^{3/4}$.
\rlabel{Исправление.} Начните с прямой подстановки известных данных, разделите на ненулевое давление и только после этого извлекайте степень.
\rlabel{Полное правильное решение.}
\[
(2\cdot10^5)V^{4/3}=3{,}2\cdot10^6.
\]
Разделим обе части на $2\cdot10^5$:
\[
V^{4/3}=\frac{3{,}2\cdot10^6}{2\cdot10^5}
=1{,}6\cdot10=16.
\]
Возводим обе части в степень $\frac34$:
\[
V=16^{3/4}=(2^4)^{3/4}=2^3=8\text{ м}^3.
\]
\rlabel{Проверка.} $8^{4/3}=(\sqrt[3]{8})^4=2^4=16$. Следовательно,
\[
(2\cdot10^5)\cdot16=3{,}2\cdot10^6,
\]
что совпадает с исходным условием.
\par\noindent\finallabel\ $\boxed{8\text{ м}^3}$.
\rlabel{Задача на закрепление.} Для газа выполняется $pV^{4/3}=8{,}1\cdot10^6$. Найдите $V$ в м$^3$, если $p=10^5$ Па.

\tasktitle{11}{Как изменение объёма влияет на давление}
\rlabel{Текст задания.} В процессе с законом $pV^a=\mathrm{const}$ объём газа уменьшается в два раза. При каком наименьшем положительном $a$ давление увеличится не менее чем в $32$ раза?
\rlabel{Ваше решение.} На представленных фотографиях решения этого задания нет.
\rlabel{Ваш ответ.} Не указан.
\par\noindent\errorlabel\ Не составлено соотношение между начальным и конечным состояниями газа.
\rlabel{Где именно возникает недочёт.} Первый необходимый шаг --- записать постоянство $p_1V_1^a=p_2V_2^a$ и учесть $V_2=\frac{V_1}{2}$.
\rlabel{Почему нужно сравнивать отношения.} В условии не заданы конкретные начальные давление и объём, но они и не требуются: при делении двух уравнений состояния эти величины сокращаются. Остаётся зависимость коэффициента роста давления только от показателя $a$.
\par\noindent\keylabel\ При уменьшении объёма в $q$ раз давление по закону $pV^a=\mathrm{const}$ увеличивается в $q^a$ раз.
\rlabel{Исправление.} Введите индексы $1$ и $2$ для двух состояний и выразите $p_2/p_1$.
\rlabel{Полное правильное решение.} Из закона процесса
\[
p_1V_1^a=p_2V_2^a
\]
получаем
\[
\frac{p_2}{p_1}=\left(\frac{V_1}{V_2}\right)^a.
\]
Так как $V_2=V_1/2$, то
\[
\frac{p_2}{p_1}=2^a.
\]
Требуется увеличение давления не менее чем в $32$ раза:
\[
2^a\ge32=2^5.
\]
Основание $2$ больше единицы, поэтому $2^a$ возрастает с ростом $a$, и
\[
a\ge5.
\]
Наименьшее значение показателя равно $5$.
\rlabel{Проверка.} При $a=5$ получаем $p_2/p_1=2^5=32$. При меньшем положительном $a$ отношение $2^a$ строго меньше $32$.
\par\noindent\finallabel\ $\boxed{5}$.
\rlabel{Задача на закрепление.} При процессе $pV^a=\mathrm{const}$ объём уменьшается в три раза. Найдите наименьшее положительное $a$, при котором давление возрастает не менее чем в $81$ раз.

\tasktitle{12}{Минимальный угол дифракционного максимума}
\rlabel{Текст задания.} Для дифракционной решётки $d\sin\varphi=k\lambda$, где $\lambda=350$ нм, а период $d\le1400$ нм. Найдите наименьший угол наблюдения второго максимума ($k=2$).
\rlabel{Ваше решение.} Вы записали формулу $d\sin\varphi=k\lambda$, подставили $1400$ и $350$ и получили соотношение $\sin\varphi=\frac{k}{4}$. Угол не найден.
\rlabel{Ваш ответ.} Не указан.
\par\noindent\errorlabel\ Не объяснено, почему для минимального угла нужно использовать максимально допустимый период $d=1400$ нм; не подставлен номер максимума и не вычислен угол.
\rlabel{Где именно остановились.} Последний рабочий шаг при выборе $d=1400$ нм действительно даёт $\sin\varphi=k/4$. Первый недостающий шаг --- объяснить этот выбор, затем положить $k=2$ и получить $\sin\varphi=1/2$.
\rlabel{Почему выбор периода требует обоснования.} Из общей формулы следует $\sin\varphi=k\lambda/d$. При фиксированных $k$ и $\lambda$ увеличение положительного $d$ уменьшает значение синуса. В рассматриваемом диапазоне углов от $0^\circ$ до $90^\circ$ вместе с синусом уменьшается и сам угол. Поэтому минимум достигается при наибольшем разрешённом $d$.
\par\noindent\keylabel\ В задачах на экстремальное значение величины сначала установите направление её изменения при изменении параметра, а затем используйте подходящую границу диапазона.
\rlabel{Исправление от последнего правильного шага.} При $d=1400$ нм
\[
\sin\varphi=\frac{k\lambda}{d}=
\frac{2\cdot350}{1400}=\frac12.
\]
\rlabel{Полное правильное решение.} Для второго максимума $k=2$, поэтому
\[
\sin\varphi=\frac{700}{d}.
\]
Чтобы угол был минимален, выбираем $d=1400$ нм. Тогда
\[
\sin\varphi_{\min}=\frac{700}{1400}=\frac12,
\qquad \varphi_{\min}=30^\circ.
\]
\rlabel{Проверка.} Подставим $\varphi=30^\circ$ и $d=1400$ нм:
\[
1400\sin30^\circ=1400\cdot\frac12=700=2\cdot350.
\]
Уравнение выполнено. При меньшем $d$ синус и угол были бы больше.
\par\noindent\finallabel\ $\boxed{30^\circ}$.
\rlabel{Задача на закрепление.} Для решётки с периодом $d\le600$ нм длина волны равна $300$ нм. Найдите наименьший угол, при котором возможен второй дифракционный максимум.

\tasktitle{13}{Параметры показательной функции по графику}
\rlabel{Текст задания.} На графике изображена функция $f(x)=a^x+b$. Найдите $f(-5)$. Из графика читаются точки $(0;-3)$ и $(-2;0)$.
\rlabel{Ваше решение.} Вы определили обе точки и верно записали
\[
-3=a^0+b=1+b,
\qquad 0=a^{-2}+b.
\]
Затем указали $b=-3$, получили $a^{-2}=3$ и $a=\pm\sqrt{1/3}$. Значение $f(-5)$ не вычислено.
\rlabel{Ваш ответ.} Не указан.
\par\noindent\errorlabel\ При нахождении $b$ потеряно слагаемое $1$. Позднее неверно добавлен отрицательный вариант основания показательной функции.
\rlabel{Где именно возникает ошибка.} Последний правильный шаг:
\[
-3=1+b.
\]
Первый неправильный шаг: $b=-3$. Из равенства следует $b=-3-1=-4$. При этом для функции $a^x$ с действительными показателями должно выполняться $a>0$ и $a\ne1$.
\rlabel{Почему это неверно.} Значение $a^0$ всегда равно $1$ при допустимом основании. Поэтому в точке с абсциссой ноль ордината равна $1+b$, а не $b$. Если подставить ошибочное $b=-3$, получается $f(0)=1-3=-2$, что не совпадает с отмеченной на графике ординатой $-3$.
\par\noindent\keylabel\ По точке на оси $Oy$ сразу найдите $b$, а по ещё одной точке определите положительное основание $a$.
\rlabel{Исправление от последнего правильного шага.}
\[
-3=1+b\ \Longrightarrow\ b=-4.
\]
По точке $(-2;0)$:
\[
0=a^{-2}-4\ \Longrightarrow\ a^{-2}=4.
\]
\rlabel{Полное правильное решение.} Получаем
\[
\frac1{a^2}=4\ \Longrightarrow\ a^2=\frac14.
\]
Так как $a>0$, выбираем только $a=\frac12$. Следовательно,
\[
f(x)=\left(\frac12\right)^x-4.
\]
Тогда
\[
f(-5)=\left(\frac12\right)^{-5}-4=2^5-4=32-4=28.
\]
\rlabel{Проверка.} $f(0)=1-4=-3$ и $f(-2)=4-4=0$. Дополнительная отмеченная на графике точка $(-3;4)$ тоже подходит: $f(-3)=8-4=4$.
\par\noindent\finallabel\ $\boxed{28}$.
\rlabel{Задача на закрепление.} График функции $g(x)=c^x+d$ проходит через точки $(0;-2)$ и $(-1;0)$. Найдите $g(-3)$.

\tasktitle{14}{Свободный член функции тангенса}
\rlabel{Текст задания.} На графике изображена функция $f(x)=a\tg x+b$. Найдите $b$.
\rlabel{Ваше решение.} На представленных фотографиях решение задания №\,14 отсутствует.
\rlabel{Ваш ответ.} Не указан.
\par\noindent\errorlabel\ Не использовано пересечение графика с осью ординат. Для определения $b$ не нужно вычислять параметр $a$.
\rlabel{Где именно возникает недочёт.} Последнего правильного шага в работе нет. Первый необходимый шаг --- подставить $x=0$ в формулу функции и прочитать по графику значение $f(0)$.
\rlabel{Почему это работает.} В точке пересечения с осью $Oy$ абсцисса равна нулю. Так как $\tg0=0$, первое слагаемое $a\tg0$ исчезает при любом допустимом $a$. Следовательно, ордината пересечения равна именно $b$.
\par\noindent\keylabel\ Для функции вида $a\tg x+b$ свободный член можно считывать непосредственно по пересечению графика с осью $Oy$.
\rlabel{Исправление.} Рассмотрите точку графика с абсциссой $0$: по рисунку её ордината равна $1$.
\rlabel{Полное правильное решение.} Из формулы
\[
f(x)=a\tg x+b
\]
при $x=0$ получаем
\[
f(0)=a\tg0+b=a\cdot0+b=b.
\]
По графику $f(0)=1$. Значит,
\[
b=1.
\]
\rlabel{Проверка.} При $b=1$ функция $f(x)=a\tg x+1$ действительно проходит через точку $(0;1)$ для любого допустимого коэффициента $a$.
\par\noindent\finallabel\ $\boxed{1}$.
\rlabel{Задача на закрепление.} График функции $h(x)=c\tg x+d$ пересекает ось $Oy$ в точке $(0;-2)$. Найдите $d$.

\vfill
\begin{center}\small\color{ink}Лёвин Артём Александрович эксклюзивно для Ярослава Гаврилова\end{center}
\end{document}